% Chapter 1 -- Rubric: PO1 (1.1), PO2 (problem identification), PO3 (3.2 objectives)
\chapter{Problem Identification and Formulation}
\label{ch:problem}

\section{Background and Motivation}
\label{sec:background}

Graphical user interfaces are the surface through which almost all consumer
software is used, and they are the surface on which almost all of its defects
become visible to a user. Verifying them is therefore unavoidable, and it is
expensive. A modern mobile application ships on a fortnightly or monthly
cadence, across several device profiles (phone, tablet, watch, television) and
several platform versions, and every release re-opens the question of whether
the parts that used to work still work.

Industry answers this question in three broad ways, and each of them is
unsatisfactory in a different direction.

\textbf{Manual exploratory testing} is what skilled human testers do: they open
the application, form a hypothesis about where it is likely to be weak, probe
it, and follow whatever they find. It is by a wide margin the most effective
technique at finding defects that nobody anticipated, because the tester is
reasoning about the product rather than replaying a script. Its weaknesses are
economic and organisational rather than technical. It does not scale with the
number of device--platform--profile combinations, the cost is linear in the
number of releases, and, most damagingly, the knowledge the tester builds up
about where the application is fragile lives in that person's head. When they
move to another project, the organisation starts again.

\textbf{Scripted automation} (Appium, Espresso, Selenium, and the page-object
patterns built on them) converts a tester's knowledge into code, which makes it
repeatable and cheap to re-run. The cost moves from execution to maintenance.
Every script encodes a set of assumptions about element identifiers, screen
layout, and navigation order, and a user-interface redesign invalidates all of
them at once. More fundamentally, a script can only check what its author
thought to check: scripted automation is a regression technique, not a
discovery technique. It confirms known behaviour; it does not explore.

\textbf{Random and search-based exploration} (Android's Monkey~\cite{monkey},
and the academic line of work on search-based GUI exploration) removes the
human from the loop entirely and drives the application with generated input.
It is cheap and it scales, but it has no test oracle worth the name. It can
detect a crash or an unhandled exception, because those are self-evident
failures, but it cannot detect that an application accepted an invalid price,
routed an order to the wrong seller, or displayed a total that does not match
its line items. Those are the defects that matter commercially, and they are
invisible without a notion of what the application was supposed to do.

\subsection{The Shift the Industry Partner Asked For}
\label{subsec:shift}

The problem addressed here was posed by an industry partner, Samsung Research
and Development Institute Bangladesh, whose brief is what fixes the project's
difficulty. The partner supplied the problem, a technical requirements document
and an early review; the design, implementation and evaluation reported in this
document are the team's own work. That review asked for one specific change: to
move from \emph{directed functional testing}, in which an agent is handed a
test objective and must carry it out, to \emph{exploratory testing}, in which
the agent must decide for itself what is worth testing next. The technical
requirements document that followed, ETA-TR-2026-001~\cite{etatr}, formalises
that shift as the evolution of a static-knowledge test generator into ``an
autonomously learning, self-optimising exploratory test engineer'', and states
the objective in one sentence that this project is ultimately evaluated
against: \emph{every execution cycle must make the agent smarter than the
last}.

The distinction is easy to state and easy to underestimate. In directed testing
the objective is an \emph{input}: a person, or an upstream tool, has already
decided what to check and usually what the correct answer is, and the agent's
job is execution. In exploratory testing the objective is an \emph{output}: the
agent must generate the objective, derive its own notion of correct behaviour,
execute, judge the result, and then let what it learned change the next
objective it generates. Exploratory testing is characterised by Kaner, Bach and
Pettichord as the \emph{simultaneous} performance of learning, test design and
test execution~\cite{kaner}; automating it means automating all three at once,
and automating the feedback between them.

\subsection{Why Exploratory Testing Is Hard}
\label{subsec:whyhard}

Six difficulties follow from that definition, and together they account for
most of the engineering effort reported in this document.

\begin{enumerate}[label=\textbf{D\arabic*.},leftmargin=3.0em]
  \item \textbf{The objective must be generated, and the search space is
    unbounded.} For any non-trivial application the set of possible action
    sequences is effectively infinite, and there is no ground truth for which
    of them is worth performing. A directed agent searches for a path to a
    given goal; an exploratory agent must first search the space of goals.
  \item \textbf{There is no stopping criterion.} A scripted suite is finished
    when its last case has run. Exploration ends when the budget ends, so every
    reported result is conditioned on a budget and no result is ever
    ``complete''. Coverage becomes something that is reported rather than
    something that is achieved.
  \item \textbf{The oracle must be derived rather than supplied.} A directed
    test case arrives with its expected result attached. An exploratory agent
    must work out what correct behaviour is from whatever specification exists and, critically, must abstain when the specification is silent rather than
    inventing an expectation. The alternative failure modes are a fabricated
    defect report and a silent pass, and \Cref{subsec:agreement} shows both are
    documented in the literature.
  \item \textbf{Credit assignment is ambiguous.} When a scripted test fails, the
    script is fixed and the application is the variable. When an exploratory run
    ends without a pass, three explanations compete: the application is wrong,
    the agent could not drive the application, or the environment blocked the
    attempt. Only the first is a defect, and conflating them inflates the reported defect count, which is the most expensive error this kind of system can make.
  \item \textbf{Non-reproducibility is intrinsic, not incidental.} Two
    exploratory runs are \emph{supposed} to do different things; that is where
    the value comes from. A regression therefore cannot be demonstrated by
    re-running the suite, and every claim about the system must be made over a
    distribution of runs rather than a single run.
  \item \textbf{Without memory, exploration is worthless.} A stateless
    exploratory agent re-treads the same ground every session. The whole value
    of the technique depends on what it learned persisting and changing what it does next, which is precisely why the partner's requirements are dominated
    by mechanisms for \emph{remembering} (\Cref{subsec:fr}).
\end{enumerate}

The literature confirms that this is the harder problem rather than merely a
different one. Of the seven agent systems reviewed in \Cref{ch:research}, five
are handed their objective: a task string, a test case, or a natural-language
test requirement~\cite{appagent,appagentx,autodroid,xuatcopilot,auitestagent}.
Only DroidAgent generates its own testing goals~\cite{droidagent}, and its
authors describe it as the first Android GUI testing technique able to do so.
The measured gap is instructive: objective-driven systems report 77\% to
88.55\% task completion~\cite{auitestagent,xuatcopilot}, while DroidAgent
completes 59\% of the tasks it set itself, with a further 15\% judged not
viable in the first place. Those figures come from different applications and
benchmarks and are not a controlled comparison, but the direction is consistent
and unsurprising: an agent that must invent its own objective will set some
objectives that cannot be achieved.

\subsection{The Gap}

The gap this project addresses is therefore specific: an automated technique
that has the \emph{oracle} of manual exploratory testing, meaning a model of
intended behaviour drawn from the specification and from what has gone wrong
before, together with the \emph{scale and repeatability} of automation, and
that, unlike both, \emph{accumulates} what it learns so that the hundredth
execution is better informed than the first.

The stakeholders affected are the partner's quality assurance engineers, whose
manual effort the system is intended to displace; application developers, who
receive its defect reports and pay the cost of false positives; and,
indirectly, end users of the applications under test.

\section{Problem Statement}
\label{sec:problem-statement}

\subsection*{Statement}

Given (i) whatever partial specification of an application exists (a software
requirements specification, a user-interface design export, a history of
previously filed defects, or none of these) and (ii) an executable instance of
that application on a real device or in a browser, construct an autonomous
agent that \emph{chooses} what to test next, executes its choice, judges the
outcome against a derived oracle, and persists what it learned, such that the
knowledge produced by each execution measurably improves the test cases
proposed in subsequent executions. No test objective is supplied from outside
the loop.

\subsection*{Engineering formulation}

Let $A$ be the application under test and let $K_t$ be the agent's knowledge
state at round $t$. The system is a loop
\[
  K_{t+1} \;=\; U\big(K_t,\; E(A,\; P(K_t))\big),
\]
where $P$ is a \emph{planner} that maps a knowledge state to a concrete test
case, $E$ is an \emph{executor} that drives $A$ through that test case and
returns an observed trajectory, and $U$ is an \emph{investigator} that reduces
a trajectory to durable, deduplicated knowledge and merges it into the graph.
The three functions are realised by separate large-language-model roles and
communicate only through $K$, which is a persistent property graph. No
conversational state is carried between calls.

The exploratory character of the problem is carried entirely by $P$. A directed
testing system would take the test case as an input and need only $E$; here $P$
must select from an unbounded space using $K_t$ as its only evidence, and the
quality of the whole system is bounded by the quality of that selection.

\begin{description}[leftmargin=2.2em,style=nextline]
  \item[Inputs] A requirements specification (free text), a user-interface
    design export, a structured defect history, and the live accessibility tree
    plus screenshot of the running application. Every one of these is optional;
    the system must degrade rather than fail when a source is absent.
  \item[Outputs] Executable test cases with steps, expected results and
    requirement citations; per-execution verdicts with an explicit fault
    attribution; atomic, deduplicated findings about the application; and a
    coverage and risk report over the requirement set.
  \item[Entities] Requirements, requirement chunks, feature areas, design
    screens and elements, observed UI states and the transitions between them,
    navigation paths, defects, test cases, execution logs, error patterns,
    strategy memory, and findings, all of them nodes in one project-scoped property graph.
  \item[Quantities to be optimised] Requirement coverage per unit cost;
    autonomy, defined as the fraction of executions that terminate with a
    verdict attributable to the application rather than to the agent or the
    environment; defect-discovery rate; duplicate-proposal rate (to be
    minimised); and monetary cost per executed test.
  \item[Quantities to be guaranteed] No destructive or irreversible action on
    the application under test; no silent failure, meaning every degradation must be recorded and surfaced; and no fabricated evidence, meaning a verdict must be
    traceable to an observation.
  \item[Assumptions] The application exposes an accessibility tree (Android) or
    a DOM (web); a dedicated, disposable test account is available; the device
    or browser is reachable from the host; and a hosted language-model API is
    available within a fixed budget.
\end{description}

\subsection*{Engineering knowledge demanded}

The formulation above is not solvable by routine application of a single body
of knowledge. It requires, at minimum: graph data modelling and query design
(the knowledge state is a property graph with constraints and vector indexes);
information retrieval, including both lexical and dense-vector similarity, to
select the fraction of $K_t$ that fits in a bounded prompt; similarity and
hashing theory, applied to the non-trivial question of when two observed
screens are the same screen; distributed-systems design, since the four
processes must share state without sharing memory; the specialist body of
knowledge of software testing, in particular test-oracle theory and fault
attribution; and empirical evaluation methodology, because the central claim of
the project, that the agent improves, is a claim about a trend and can only be
established statistically.

\section{Complexity of the Problem}
\label{sec:complexity}
A complex engineering problem must have characteristic P1 and some or all of
P2--P7 (\Cref{tab:complexity}). P1 requires in-depth knowledge at the level of
one or more of the knowledge profiles K3--K8 listed in
\Cref{tab:knowledge-profile}.

\begingroup
\small
\setlength{\tabcolsep}{5pt}
\begin{xltabular}{\textwidth}{@{}l>{\hsize=1.10\hsize}Y>{\hsize=0.90\hsize}Y@{}}
  \caption{Complex engineering problem characteristics (P1--P7) and evidence from this project}
  \label{tab:complexity}\\
  \toprule
  Code & Definition & Evidence from this project \\
  \midrule
  \endfirsthead
  \caption[]{Complex engineering problem characteristics (P1--P7) \emph{(continued)}}\\
  \toprule
  Code & Definition & Evidence from this project \\
  \midrule
  \endhead
  \midrule
  \multicolumn{3}{r@{}}{\footnotesize\itshape continued on the next page}\\
  \endfoot
  \bottomrule
  \endlastfoot

  P1 & \textbf{Depth of knowledge required.} Cannot be resolved without in-depth engineering knowledge at the level of one or more of K3, K4, K5, K6, or K8, which allows a fundamentals-based, first-principles analytical approach. &
  \textbf{K3, K4, K5, K6 and K8 are all engaged.} K3: the state-identity problem was solved from first principles as a cascade of set-similarity measures (exact signature, Jaccard index, containment, structural-skeleton containment) with a perceptual image hash as the final tie-break, each with an analytically chosen threshold and guard. K4: test-oracle theory and fault attribution, taken from the software-testing body of knowledge, are what separate an application defect from an agent failure. K5: the four-process architecture and the single-writer graph discipline are design-knowledge decisions. K6: current practice in Android accessibility services, Playwright, Neo4j and hosted language-model tool-calling. K8: seven recent systems were read in full and a further seven surveyed, as reported in \Cref{ch:research}. \\
  \addlinespace
  P2 & \textbf{Range of conflicting requirements.} Involves wide-ranging or conflicting technical, engineering, and other issues. &
  Four irreducible conflicts were managed rather than eliminated: \emph{autonomy versus safety}, where the agent must act freely but must never log out, delete an account, or make a payment; \emph{coverage versus cost}, where every additional test is a metered charge against a fixed budget; \emph{generality versus grounding}, where knowledge of the specific application makes the agent effective and simultaneously makes it useless on the next application; and \emph{exploration versus reproducibility}, which is difficulty D5, where the agent's value comes from non-deterministic exploration while the report it produces must be defensible evidence. \\
  \addlinespace
  P3 & \textbf{Depth of analysis required.} Has no obvious solution and requires abstract thinking and originality in analysis to formulate suitable models. &
  Three sub-problems had no off-the-shelf answer. (i) \emph{What should be tested next?} This is difficulty D1, and it has no ground truth; the answer adopted is to rank an unbounded space by evidence held in the knowledge graph (uncovered requirements, defect-dense areas, regression risk) rather than by any fixed heuristic. (ii) \emph{When are two screens the same screen?} A naive signature fragmented one wizard into nine distinct states, destroying the navigation model; the multi-stage cascade described under P1 was designed in response. (iii) \emph{Whose fault was this failure?} Difficulty D4: an execution that ends without a pass is ambiguous between a defect, an agent limitation and an environmental block, and a three-way attribution taxonomy with a single authoritative definition was introduced to prevent conflation. \\
  \addlinespace
  P4 & \textbf{Familiarity of issues.} Involves infrequently encountered issues. &
  The system embeds non-deterministic components inside a harness whose own output is a correctness judgement, an inversion of the usual assumption that a test harness is deterministic. Further non-routine issues: applications built with cross-platform toolkits expose no stable resource identifiers and report a single activity on every screen, defeating conventional element location; and prompt injection becomes a live concern because the application's own on-screen text is fed to a model that also decides what to do next. \\
  \addlinespace
  P5 & \textbf{Extent of applicable codes.} Is outside problems encompassed by standards and codes of practice for professional engineering. &
  ISO/IEC/IEEE 29119~\cite{iso29119} gives a vocabulary and a process for testing but presumes that test cases are authored by people and that test design is deterministic. It has nothing to say about an agent that authors its own test cases at run time, nor about how to report coverage when the test set is not fixed in advance, which is exactly the regime difficulty D2 describes. No standard governs the correctness of a language-model-produced verdict. The applicable codes therefore had to be supplemented by project-specific rules, documented in \Cref{sec:standards}. \\
  \addlinespace
  P6 & \textbf{Extent of stakeholder involvement and conflicting requirements.} Involves diverse groups of stakeholders with widely varying needs. &
  The industry partner wants learning that transfers across device profiles and platforms; a quality-assurance lead wants runs that can be reproduced and audited; application developers want a near-zero false-positive rate, because each false defect report costs them investigation time; the academic supervisor requires a defensible empirical evaluation; and the end users of the application under test require that testing never damage live data. These pull in different directions, most sharply transfer-driven generalisation against reproducibility. \\
  \addlinespace
  P7 & \textbf{Interdependence.} Is a high-level problem including many component parts or sub-problems. &
  The partner's requirements document itself decomposes the work into eight requirement families (ETA-REQ-301 to 308) with an explicit dependency ordering, and the delivered system comprises four cooperating processes, two execution platforms (Android and web), three distinct language-model roles, nine pluggable knowledge sources and a single shared graph schema. A change to the attribution taxonomy, for example, propagates to the executor, the investigator, the risk score, the coverage report and the generated PDF. \\
\end{xltabular}
\endgroup

\begin{table}[htbp]
  \centering
  \caption{Knowledge profiles referenced by P1}
  \label{tab:knowledge-profile}
  \begin{tabularx}{\textwidth}{@{}lY@{}}
    \toprule
    Code & Description \\
    \midrule
    K3 & A systematic, theory-based formulation of engineering fundamentals required in the engineering discipline. \\
    K4 & Engineering specialist knowledge that provides theoretical frameworks and bodies of knowledge for the accepted practice areas in the engineering discipline; much is at the forefront of the discipline. \\
    K5 & Knowledge that supports engineering design in a practice area. \\
    K6 & Knowledge of engineering practice (technology) in the practice areas in the engineering discipline. \\
    K8 & Engagement with selected knowledge in the research literature of the discipline. \\
    \bottomrule
  \end{tabularx}
\end{table}

\section{Project Objectives and Scope}
\label{sec:objectives}

\subsection*{Objectives}

The objectives below are stated so that each one can be checked;
\Cref{ch:evaluation} evaluates the delivered system against them. O2 and O4
together are the exploratory requirement of \Cref{subsec:shift}: the agent must
choose its own objectives, and what it learns must change the choices it makes
afterwards.

\begin{enumerate}[label=\textbf{O\arabic*.},leftmargin=3.2em]
  \item \textbf{Build a persistent, project-scoped knowledge graph} that ingests
    a requirements specification, a user-interface design export and a defect
    history, and that is additionally grown from the agent's own observations of
    the running application.
    \emph{Verified by:} ingesting each source independently and querying the
    resulting nodes and relationships.

  \item \textbf{Select and generate executable, non-duplicate test cases from
    that graph without an externally supplied objective}, each citing the
    requirement it covers.
    \emph{Verified by:} the proportion of generated test cases that carry a
    valid requirement citation, the rejection rate of re-worded duplicates, and
    the absence of any human-supplied objective in the loop.

  \item \textbf{Execute generated test cases autonomously} against a real
    Android device and against a live web application, and return a verdict with
    an explicit three-way fault attribution.
    \emph{Verified by:} the autonomy rate, meaning executions that terminate with a verdict about the application rather than about the agent, measured over a campaign.

  \item \textbf{Close the learning loop}, so that execution outcomes are
    persisted as durable findings, navigation paths, error patterns and strategy
    scores, and demonstrably alter the next round's proposals.
    \emph{Verified by:} showing that a proposal changes when, and only when, the
    corresponding knowledge is present in the graph.

  \item \textbf{Satisfy the eight requirement families of
    ETA-TR-2026-001}~\cite{etatr} (ETA-REQ-301 to 308) and its two non-functional
    requirements.
    \emph{Verified by:} the requirement-to-implementation traceability matrix in
    \Cref{sec:requirements}, and the partner's stated acceptance criteria.

  \item \textbf{Keep the system application-agnostic.} No knowledge of any
    particular application may be embedded in code or prompts; all
    application-specific information must arrive through configuration.
    \emph{Verified by:} an automated generality audit over the source tree, and
    by running the same unmodified system against more than one application.

  \item \textbf{Operate within a fixed monetary budget}, with the cost per
    executed test measured and reported rather than estimated.
    \emph{Verified by:} reconciling the provider's billing record against the
    number of executions.

  \item \textbf{Provide an operator surface}: a live dashboard and an exportable report, so that a run can be watched, audited and handed to
    someone who was not present for it.
    \emph{Verified by:} demonstration against a live campaign.
\end{enumerate}

\subsection*{Scope}

\textbf{Within scope.} Autonomous exploratory testing of (i) Android
applications on a physical device or emulator, driven through the platform
accessibility service, and (ii) web applications, driven through a headless or
headed browser; ingestion of requirements, design exports and defect histories;
a persistent knowledge graph with learning across runs; fault attribution and
self-healing retry; risk-weighted prioritisation; and an operator dashboard
with report export.

\textbf{Explicitly out of scope.} Unit, integration, performance, load and
security testing of the application under test; this project addresses
system-level testing through the user interface only. Directed execution of a
human-authored test suite: the system may be given a specification, never an
individual objective, since supplying objectives would remove the exploratory
problem the partner asked to have solved. Testing of flows that cannot be
completed without an irreversible side effect, specifically account
registration, one-time-password verification, payment and account deletion;
these are guarded against rather than exercised, and the corresponding
requirements are marked as preconditions already satisfied. Automated repair of
discovered defects. Training or fine-tuning of language models: the project
uses hosted models through an API and treats them as components. Formal
certification of the agent's verdicts; a reported defect is a candidate for
human confirmation, not a confirmed defect.

\textbf{Deferred.} Cross-device-profile transfer beyond the \texttt{application}
axis, and the seeded-defect experiment needed to measure precision and recall
against a known ground truth, are specified and partially implemented but not
validated within this session; both are carried into \Cref{ch:conclusion} as
future work.
